\section{Assessing Ba$_5$Y$_{12}$Zn[O(SiO$_4$)]$_8$ as a Structural Analogue}

Our search of OpenAlex, Crossref, and the accessible full-text corpus returned no primary crystal-structure entry for Ba$_5$Y$_{12}$Zn[O(SiO$_4$)]$_8$. The ``yttrium barium silicates'' that did surface are mostly luminescent thin films and oxyapatite-type materials whose stoichiometry, space group, and Zn occupancy all differ from the target formula \cite{gupta2021combustion,yang2021cyan}. The defensible label at present is therefore ``candidate structural analogue awaiting verification'' rather than ``known isostructural compound''.

A tiered comparison is still possible on that footing. BaZn$_2$Si$_2$O$_7$ gives a direct reference for the corner-sharing ZnO$_4$--SiO$_4$ network and the Ba polyhedra \cite{lin1999phase,available2020materials,kaiser2002crystal}; the yttrium barium silicate oxyapatite literature gives a compositional reference in which Y\slash Ba coexist with isolated silicate groups, but it does not show that the two topologies are the same \cite{gupta2021combustion,yang2021cyan,wei2013a}. If powder or single-crystal data for the target phase become available, the degree of silicate polymerization, the Zn coordination number, the connectivity of the Y polyhedra, the oxygen-vacancy positions, and the space group should be checked in that order, rather than simply comparing element ratios in the formula.

The high Y content of Ba$_5$Y$_{12}$Zn[O(SiO$_4$)]$_8$ also raises a charge-compensation issue. Formal valence counting calls for additional O$^{2-}$ or for mixed cation valence, so any synthesis attempt must confirm the oxygen stoichiometry and valence states by thermogravimetric analysis, XPS\slash XANES, or neutron diffraction. The thermally driven transition of BaZn$_2$Si$_2$O$_7$ can guide the choice of heat-treatment window but cannot be extrapolated directly to the Y-rich system. Table~\ref{tab:analogue} lists the grounds for comparison by evidence tier.

\begin{table}[bp]
\centering
\caption{Evidence tiers for the structural-analogue comparison.}
\label{tab:analogue}
\begin{tabular}{P{\dimexpr 0.1835\textwidth-2\tabcolsep\relax}P{\dimexpr 0.3538\textwidth-2\tabcolsep\relax}P{\dimexpr 0.4628\textwidth-2\tabcolsep\relax}}
\toprule
    \thead{Comparison tier} & \thead{Evidence for BaZn$_2$Si$_2$O$_7$} & \thead{Implication for Ba$_5$Y$_{12}$Zn[O(SiO$_4$)]$_8$}\\
\midrule
    Silicate polymerization & Corner-sharing ZnO$_4$\slash SiO$_4$ network \cite{lin1999phase} & A topological starting point to be checked; not proof of isostructurality\\
    \addlinespace[3pt]
    Cation coordination & Five-coordinate Ba, four-coordinate Zn \cite{available2020materials} & Y and Zn coordination must be measured, not inferred from ionic radii\\
    \addlinespace[3pt]
    Thermal stability & Transition at about 280--300~$^\circ$C \cite{thieme2015ba1} & A reference for the thermal-analysis program only\\
\bottomrule
\end{tabular}
\end{table}
